\section{Default Winner Architecture}

Once the standard identifies a default, continuing to optimize for every runner-up becomes noise. The agent needs one default map. The audit needs one default policy. CI needs one default set of gates. Exceptions remain available, but they must be named and versioned.

\begin{figure*}[t]
\centering
\begin{minipage}[t]{0.48\textwidth}
\begin{lstlisting}[style=codebox]
repo/
  apps/
    web/              TS/React/Vite UI
    api/              Rust HTTP/RPC edge
  crates/
    domain/           pure invariants
    application/      use cases/authz/tx
    adapters/         Postgres/queues/APIs
    workers/          jobs/replay/backoff
  contracts/          OpenAPI/Protobuf/JSON
  db/                 migrations/constraints
  python/ai-service/  model/eval boundary
  ops/                CI/security/telemetry
\end{lstlisting}
\end{minipage}
\hfill
\begin{minipage}[t]{0.48\textwidth}
\begin{lstlisting}[style=codebox]
repo/
  AGENTS.md           short root router
  agent/
    owner-map.json    path -> owner
    test-map.json     path -> proof lane
    proof-lanes.toml  runnable commands
    generated-zones.toml
    standard-version.toml
    agent/repo-score.json
  docs/
    decisions/
    exceptions/
    repair/
    runbooks/
\end{lstlisting}
\end{minipage}
\caption{Ownership spine and agent-control spine for the default stack.}
\label{fig:repo-tree-main}
\end{figure*}

\begin{table*}[t]
\caption{Default ownership cells for the winning stack.}
\label{tab:ownership}
\centering
\small
\begin{tabularx}{\textwidth}{L{0.18\textwidth} Y Y}
\toprule
\textbf{Cell} & \textbf{Owns} & \textbf{Must not own} \\
\midrule
\texttt{apps/web} & React UI, route state, forms, local validation, generated API clients, Playwright tests. & Secrets, durable truth, direct SQL, core authorization, handwritten API mirrors. \\
\texttt{apps/api} & Rust HTTP/RPC edge, request normalization, auth middleware, idempotency keys, response shaping. & Domain invariants hidden in handlers, scattered SQL, UI concerns. \\
\texttt{domain} crate & IDs, value objects, invariants, state machines, pure decisions, stable error types. & I/O, env reads, network, filesystem, database clients, framework code. \\
\texttt{application} crate & Commands, authz, transaction orchestration, idempotency, workflow use cases. & Transport details, SQL drivers, UI state, prompt/model logic. \\
\texttt{adapters} crate & PostgreSQL access, queue/streaming clients, external APIs, filesystem, environment, provider clients. & Domain rules, authorization policy, or event schema ownership. \\
\texttt{workers} crate & Async jobs, retries, backoff, workflow replay, scheduled tasks. & Duplicated invariants or unowned product truth. \\
\texttt{contracts} & OpenAPI, Protobuf, JSON Schema, generated stubs and clients. & Hand-edited generated outputs or business logic. \\
\texttt{db} & Migrations, constraints, indexes, RLS, seed rules, schema snapshots. & Application orchestration or hidden business process. \\
\texttt{python/ai-service} & Inference, embeddings, evals, prompt experiments, offline transforms, typed AI API. & Product truth, authz, billing, direct production DB ownership. \\
\texttt{ops} & CI, release, telemetry, secret scanning, SBOM/SCA, provenance, policy gates. & Feature behavior hidden behind manual operations. \\
\bottomrule
\end{tabularx}
\end{table*}

The filesystem should be a routing table. A patch touching a payment invariant should land in Rust domain/application and PostgreSQL constraints. A patch touching a form should land in TypeScript UI and generated client contracts. A patch touching model behavior should land in the bounded Python service and contract tests. If the agent must infer ownership from convention, the repo has already failed.

Streaming follows the same rule. Event schemas belong in \texttt{contracts/}; generated event types belong in declared generated zones; Kafka, Tansu, Iggy, Fluvio, NATS, Redis Streams, or equivalent clients belong in \texttt{crates/adapters/queues} or an explicitly declared adapter path. Domain and application code may name business events, but it should not own broker protocol details.
